Rozhodca v~zápasení komunikuje svoje rozhodnutia pomocou
gest, ktoré následne ručne zaznamenáva do bodovacej karty,
čo zdržuje priebeh zápasu a~je náchylné na chyby. Práca
skúma, či je možné tento proces zjednodušiť pomocou
bežne dostupných inteligentných hodiniek bez nutnosti
metód strojového učenia. Navrhnuté riešenie využíva
hodinky Amazfit T-Rex~3 so~Zepp~OS nosené na ľavom
zápästí rozhodcu, ktoré snímajú pohyb pomocou akcelerometra
a~gyroskopu. Dáta sa v~reálnom čase vyhodnocujú v~sprievodnej
Android aplikácii pomocou prahových hodnôt nakalibrovaných
pre osem signálov rozhodcu podľa pravidiel \Gls{uww}. Pri
návrhu sme vychádzali z~prirodzenej gestikulácie rozhodcu,
aby si nemusel osvojovať umelé pohyby. Spoľahlivosť systému
sme overili v~troch testovacích fázach, ktoré sa postupne
približovali k~reálnym podmienkam použitia -- od izolovaných
gest cez kompletný zápasnícky scenár až po testovanie
s~externými používateľmi bez predchádzajúcej skúsenosti.
Hlavným zisteným obmedzením je závislosť prahového algoritmu
na individuálnom štýle vykonávania gest, ktorú možno do
budúcna riešiť kalibračnou fázou prispôsobenou každému
rozhodcovi.